\documentclass[11pt]{article}
\usepackage[a4paper,margin=27mm]{geometry}
\usepackage[T1]{fontenc}
\usepackage{lmodern}
\usepackage{microtype}
\usepackage{amsmath,amssymb,amsthm,mathtools}
\usepackage{enumitem}
\usepackage{xcolor}
\usepackage[colorlinks=true,linkcolor=blue!55!black,citecolor=blue!55!black,urlcolor=blue!55!black]{hyperref}
\hypersetup{
  pdftitle={Characteristic-Free Bounds and Sharp Gauss-Fiber Examples for Point-Span Tangent Absorption},
  pdfauthor={Lluis Eriksson},
  pdfsubject={Algebraic geometry; Gauss maps; first jets; interpolation},
  pdfkeywords={Gauss map, tangent absorption, first jets, fat points, positive characteristic}
}

\newtheorem{theorem}{Theorem}[section]
\newtheorem{proposition}[theorem]{Proposition}
\newtheorem{lemma}[theorem]{Lemma}
\newtheorem{corollary}[theorem]{Corollary}
\newtheorem{remark}[theorem]{Remark}
\newtheorem{definition}[theorem]{Definition}
\newcommand{\PP}{\mathbf P}
\newcommand{\CC}{\mathbf C}
\newcommand{\kk}{\mathbf k}
\newcommand{\cO}{\mathcal O}
\newcommand{\Span}{\operatorname{Span}}
\newcommand{\Sym}{\operatorname{Sym}}
\newcommand{\Gr}{\operatorname{Gr}}

\title{Characteristic-Free Bounds and Sharp Gauss-Fiber Examples\\
for Point-Span Tangent Absorption}
\author{Lluis Eriksson\\\small Independent researcher}
\date{August 24, 2026}

\begin{document}
\maketitle

\begin{abstract}
Let $X^d\subset\PP(V)$ be a smooth projective variety over an algebraically
closed field, let $H=\cO_X(1)$, and use the complete embedding defined by
$H^m$, $m\geq3$.  A finite reduced set $Z\subset X$ is point-span
tangent-absorbing if the span $S_Z$ of its evaluation functionals contains the
affine embedded tangent space at every support.  We prove, in arbitrary
characteristic,
\[
 \dim S_Z,\ |Z|\geq
 \max\left\{J(d,m),
 \min\left((d+1)(d+2),\binom{d+m}{d}\right)\right\},
\]
where
\[
 J(d,m)=\left\lfloor\frac{m+1}{2}\right\rfloor(d+1)
       +\begin{cases}1,&m\text{ even},\\0,&m\text{ odd}.
         \end{cases}
\]
The proof retains the escape-or-Gauss-fiber alternative.  The
characteristic-free step is a minimal-degree argument: in characteristic
$p$, a polynomial with all first partials zero is a $p$th power, and
radicality replaces the characteristic-zero derivative contradiction.

We give two complementary geometric results.  If the original embedding is
the complete $H$-embedding and $H=A\otimes B$ with both factors separating
distinct points, its Gauss map is injective and only the strong branch can
occur.  Conversely, over
$\CC$, for every $d\geq2$ and $m\geq3$ we construct a smooth hypersurface
whose exceptional Gauss branch contains an absorbing set with
\[
 \dim S_Z=|Z|=\binom{d+m}{d}.
\]
Thus the binomial estimate is exact within the exceptional branch in all
these dimensions and degrees.  The construction is existential and uses
finite jet interpolation and Bertini; no classification or broad priority
claim is made.  A companion paper proves the sharper exact global binomial
floor in characteristic zero; the present paper's distinct contribution is
the characteristic-free alternative and its accompanying Gauss-map geometry.
\end{abstract}

\section{Introduction and statements}

Let $\kk$ be an algebraically closed field, let
$X^d\subset\PP(V)$ be smooth and integral, and put $H=\cO_X(1)$.  Write
\[
 \varphi_m=\varphi_{|H^m|}:X\longrightarrow
 \PP\bigl(H^0(X,H^m)^*\bigr)
\]
for the complete embedding.  For a nonempty finite reduced set $Z\subset X$,
set
\[
 S_Z:=\Span_{\kk}\{\operatorname{ev}_p:p\in Z\}
       \subset H^0(X,H^m)^*.
\]

\begin{definition}
The set $Z$ is \emph{point-span tangent-absorbing} with respect to $H^m$ if
\[
 \widehat T_{\varphi_m(p)}\varphi_m(X)\subseteq S_Z
 \qquad(p\in Z),
\]
where the left side is the affine tangent space to the cone over the complete
$H^m$-embedding.
\end{definition}

Let
\[
 \gamma_H:X\longrightarrow\Gr(d+1,V),\qquad
 p\longmapsto\widehat T_pX
\]
be the Gauss map of the original $H$-embedding.  The main result is the
following.

\begin{theorem}[Characteristic-free rank alternative]\label{thm:main}
Let $\kk$ be algebraically closed of arbitrary characteristic, let
$X^d\subset\PP(V)$ be smooth and integral with $d\geq1$, and let $m\geq3$.
If $Z\subset X$ is a nonempty finite reduced point-span tangent-absorbing set
for the complete $H^m$-embedding, then at least one of the following
nonexclusive alternatives holds:
\begin{enumerate}[label=(\alph*)]
\item $Z$ contains $d+2$ supports whose first infinitesimal neighborhoods
impose independent conditions on $H^m$, and hence
\[
 \dim S_Z,\ |Z|\geq(d+1)(d+2);
\]
\item there is a vector subspace $W\subset V$, $\dim W=d+1$, such that
\[
 Z\subset X\cap\PP(W),\qquad
 \widehat T_pX=W\quad(p\in Z),
\]
and
\[
 \dim S_Z,\ |Z|\geq\binom{d+m}{d}.
\]
\end{enumerate}
Consequently,
\begin{equation}\label{eq:two-branch}
 \dim S_Z,\ |Z|\geq
 \min\left\{(d+1)(d+2),\binom{d+m}{d}\right\}.
\end{equation}
\end{theorem}

The next bound is independent of the Gauss alternative and grows with $m$.

\begin{theorem}[Growing jet-block bound]\label{thm:growing}
Under the hypotheses of Theorem~\ref{thm:main}, put
\[
 J(d,m):=\left\lfloor\frac{m+1}{2}\right\rfloor(d+1)
          +\begin{cases}1,&m\text{ even},\\0,&m\text{ odd}.
            \end{cases}
\]
Then
\[
 \dim S_Z,\ |Z|\geq J(d,m).
\]
In particular, combining this with \eqref{eq:two-branch} gives the bound in
the abstract, and $J(d,m)>(d+1)(d+2)$ for every $m\geq2d+4$.
\end{theorem}

Our first geometric corollary rules out the exceptional branch for a broad
class of original embeddings.

\begin{corollary}[Factorized polarizations]\label{cor:factorized}
Assume in addition that the original embedding is the complete $H$-embedding,
so $V=H^0(X,H)^*$, and that $H=A\otimes B$, where $A$ and $B$ each separate
every ordered pair of distinct closed points.  Then the Gauss map $\gamma_H$ is
injective on closed points, alternative~\textup{(a)} of
Theorem~\ref{thm:main} holds, and
\[
 \dim S_Z,\ |Z|\geq(d+1)(d+2).
\]
In particular this applies when $H=L^r$ for a very ample line bundle $L$ and
$r\geq2$.
\end{corollary}

\begin{remark}[Why completeness is necessary]\label{rem:incomplete-counterexample}
The factorization alone does not control the Gauss map of an arbitrary
incomplete subsystem.  Indeed
\[
 [s:t]\longmapsto
 [s^5+st^4:s^4t+t^5:s^3t^2:s^2t^3]
\]
embeds $\PP^1$ in $\PP^3$ by a subsystem of $H^0(\PP^1,\cO(5))$.
The four coordinates have no common zero: at either endpoint one of the first
two is nonzero, and on $st\ne0$ the last two are nonzero.
On $st\ne0$ the ratio of the last two coordinates recovers $t/s$; near
$[1:0]$ and $[0:1]$, respectively, the ratios of the second to the first and
the first to the second recover local parameters.  Thus the map is injective
and immersive.  Nevertheless the affine tangents at the two endpoints are
both $\Span(e_0,e_1)$.  Here
$\cO(5)=\cO(2)\otimes\cO(3)$ with both factors very ample.  Hence the complete
embedding hypothesis in Corollary~\ref{cor:factorized} cannot simply be
dropped.
\end{remark}

The opposite behavior also occurs.  The following statement makes exact only
the estimate inside branch~\textup{(b)}; it does not say that the binomial is
the global minimum when the other branch has the smaller bound.

\begin{theorem}[Sharp exceptional Gauss branch]\label{thm:sharp-gauss}
For every $d\geq2$ and $m\geq3$, set
\[
 N=\binom{d+m}{d},\qquad e=3N.
\]
There exist a smooth integral complex hypersurface
$X^d\subset\PP^{d+1}$ of degree $e$, a hyperplane
$\Lambda\cong\PP^d$, and a reduced set $Z\subset X\cap\Lambda$ of $N$
points such that
\[
 \widehat T_pX=\widehat\Lambda\quad(p\in Z)
\]
and $Z$ is point-span tangent-absorbing for the complete
$\cO_X(m)$-embedding, with
\[
 \dim_{\CC}S_Z=|Z|=N.
\]
Thus the binomial lower bound in alternative~\textup{(b)} of
Theorem~\ref{thm:main} is attained for every such $(d,m)$.
\end{theorem}

The proofs are self-contained apart from standard Bertini and cohomology of
projective space.  The finite-jet constructions are related to jet ampleness
\cite{BeltramettiDiRoccoSommese1999}; the rank defect sits at the extreme end
of Terracini-type interpolation phenomena
\cite{BallicoChiantini2021,GaluppiEtAl2023}.  Earlier public working-note
forms of the hypersurface and factorization arguments appear in
\cite{ErikssonB219,ErikssonB220}.  Those notes are antecedents, not external
validation.  The companion ARR paper \cite{ErikssonExactARR2026} proves the
sharper exact universal binomial floor in characteristic zero by an adapted
finite projection and a symbolic-power argument.  The present paper neither
supersedes nor reproves that sharper characteristic-zero statement: it
extends the rank framework to arbitrary characteristic, adds the growing
bound $J(d,m)$, and isolates complementary Gauss-map criteria and examples.
No exhaustive novelty or priority assertion is made.

\section{First jets and the annihilator criterion}

For $p\in X$, let $2p$ be the first infinitesimal neighborhood defined by
$\mathfrak m_p^2$.  It has length $d+1$.  Restriction to $2p$ is surjective
for a very ample line bundle, and the dual of its target is the affine tangent
space to the cone over the associated embedding.

Let $2Z=\coprod_{p\in Z}2p$.  The point span and the tangent span have
annihilators
\[
 S_Z^\perp=H^0(X,\mathcal I_Z\otimes H^m),
\]
and
\[
 \left(\sum_{p\in Z}
 \widehat T_{\varphi_m(p)}\varphi_m(X)\right)^\perp
 =H^0(X,\mathcal I_{2Z}\otimes H^m).
\]
Since each evaluation line lies in its affine tangent space, double
annihilators give the following exact criterion.

\begin{lemma}[Kernel criterion]\label{lem:kernel}
The set $Z$ is point-span tangent-absorbing if and only if
\[
 H^0(X,\mathcal I_Z\otimes H^m)
 =H^0(X,\mathcal I_{2Z}\otimes H^m).
\]
Equivalently, restriction to $Z$ and restriction to $2Z$ have the same rank.
\end{lemma}

This formulation is scheme-theoretic and characteristic-free.

\section{Independent first-jet blocks}

We use two kinds of finite interpolation.  The first exploits projective
independence and is efficient already in degree three.

\begin{lemma}[Projectively independent supports]\label{lem:projective-blocks}
Let $S\subset X$ be a finite set whose representatives in $V$ are linearly
independent.  If $m\geq3$, then $2S$ imposes independent conditions on $H^m$.
\end{lemma}

\begin{proof}
Order $S=\{p_1,\ldots,p_r\}$.  A single first neighborhood is separated by
$H^m$.  Suppose the first neighborhoods at
$p_1,\ldots,p_{j-1}$ have been separated.  Projective independence supplies
a linear form $a$ vanishing on their representative span and nonzero at
$p_j$.  Then $a^2$ vanishes on all the previous first neighborhoods and is a
unit on $2p_j$.  Since $H^{m-2}$ is very ample, its sections restrict
surjectively to $2p_j$; multiplication by the unit $a^2$ preserves that local
surjectivity without changing the old data.  Induction proves the claim.
\end{proof}

This direct proof is the characteristic-free specialization needed here.
Over $\CC$, the same conclusion also follows from Ballico--Chiantini's
product criterion with $L_1=H^2$, $L_2=H^{m-2}$ and their condition
$\dagger$ \cite[Definition~3.1 and Theorem~3.5]{BallicoChiantini2021}.

The second interpolation statement works for arbitrary distinct supports and
is the direct input for Theorem~\ref{thm:growing}.

\begin{lemma}[Arbitrary-support jet blocks]\label{lem:arbitrary-blocks}
Let $p_1,\ldots,p_r$ be distinct points of $X$.
\begin{enumerate}[label=(\roman*)]
\item The union $2p_1\cup\cdots\cup2p_r$ imposes independent conditions on
$H^{2r-1}$.
\item If $q$ is a further point, then
$2p_1\cup\cdots\cup2p_r\cup q$ imposes independent conditions on $H^{2r}$.
\end{enumerate}
Both conclusions remain true after increasing the exponent of $H$.
\end{lemma}

\begin{proof}
For every ordered pair of distinct supports choose a section of $H$ that
vanishes at the second support and is nonzero at the first.

For part~(i), fix $i$ and multiply the squares of separators vanishing at all
$p_j$ with $j\neq i$.  The product $P_i\in H^0(X,H^{2r-2})$ is a unit at
$p_i$ and vanishes on $2p_j$ for $j\neq i$.  Multiplying $P_i$ by sections of
$H$, whose first jets span at $p_i$, realizes arbitrary data on $2p_i$ and
zero data on every other block.  Summing over $i$ proves surjectivity.

For part~(ii), multiply the preceding $P_i$ by one separator vanishing at
$q$ and nonzero at $p_i$, then multiply by a section of $H$ to prescribe the
first jet at $p_i$.  This has total degree $2r$ and kills the simple value at
$q$.  Conversely, the product of squared separators vanishing at each
$p_i$ and nonzero at $q$ is a section of $H^{2r}$ that vanishes on all double
blocks and is a unit at $q$.  These sections separate the whole mixed scheme.

To raise the degree, multiply the constructed sections by a section of a
positive power of $H$ that is nonzero at every selected support.  Such a
section exists because $\kk$ is infinite and a finite union of proper linear
subspaces cannot exhaust the space of sections.  Multiplication by this local
unit preserves the prescribed truncated jets.
\end{proof}

Lemma~\ref{lem:arbitrary-blocks} is also a concrete instance of the tensor
product behavior of jet-ampleness established in
\cite[Proposition~2.3]{BeltramettiDiRoccoSommese1999}.

\section{Escape or concentration in a Gauss fiber}

\begin{lemma}[Tangent escape separator]\label{lem:separator}
Let $0\neq v\in W\subsetneq V$ represent $p=[v]\in X$, and let
$u\in\widehat T_pX\setminus W$.  For every $m\geq1$ there is a section of
$H^m$ that vanishes on $X\cap\PP(W)$ but has nonzero first jet at $p$ in the
direction $u$.
\end{lemma}

\begin{proof}
Choose $\lambda\in V^*$ with $\lambda|_W=0$ and $\lambda(u)\neq0$, and choose
$\mu\in V^*$ with $\mu(v)\neq0$.  The restriction to $X$ of
$f=\lambda\mu^{m-1}$ vanishes on $X\cap\PP(W)$, while
\[
 df_v(u)=\lambda(u)\mu(v)^{m-1}\neq0.
\]
No numerical factor depending on $m$ is used.
\end{proof}

\begin{lemma}[Support alternative]\label{lem:support-alternative}
Assume $Z$ is point-span tangent-absorbing.  Let
$p_1,\ldots,p_r\in Z$ have linearly independent representatives and put
$W=\Span(v_1,\ldots,v_r)$.
\begin{enumerate}[label=(\roman*)]
\item If $r\leq d$, some point of $Z$ lies outside $\PP(W)$.
\item If $r=d+1$, either some point of $Z$ lies outside $\PP(W)$, or
\[
 Z\subset\PP(W),\qquad \widehat T_pX=W\quad(p\in Z).
\]
\end{enumerate}
\end{lemma}

\begin{proof}
Suppose $Z\subset\PP(W)$.  If $r\leq d$, then
$\dim W=r<d+1=\dim\widehat T_{p_1}X$, so
Lemma~\ref{lem:separator} produces a section annihilating $S_Z$ but not the
tangent space at $p_1$, contrary to absorption.

If $r=d+1$, the spaces $W$ and $\widehat T_pX$ have the same dimension.  If
they differ for some $p\in Z$, choose a tangent direction outside $W$ and
apply the same separator.  Hence the no-escape case forces equality at every
support.
\end{proof}

\section{Unisolvence in the Gauss branch in all characteristics}

\begin{lemma}[Gauss-branch unisolvence]\label{lem:gauss-unisolvence}
Let $W$ have dimension $d+1$, let $Z\subset\PP(W)$ be a nonempty finite
reduced set, and assume
\[
 \widehat T_pX=W\qquad(p\in Z).
\]
If $Z$ is point-span tangent-absorbing for $H^m$, then, for
$R=\Sym(W^*)$ and $I=I(Z)\subset R$,
\[
 I_m=0.
\]
Consequently,
\[
 \dim S_Z,\ |Z|\geq\binom{d+m}{d}.
\]
\end{lemma}

\begin{proof}
Let $F\in I_m$.  The inclusion $W\subseteq V$ induces a surjection
$\Sym^m(V^*)\twoheadrightarrow\Sym^m(W^*)$, so $F$ extends to a degree-$m$
form on $V$.  Its restriction to $X$ vanishes on $Z$, hence pairs to zero
with $S_Z$.  Absorption forces its first jet along $X$ to vanish at every
$p\in Z$.  Since $\widehat T_pX=W$, this is the full projective first jet of
$F$ on $\PP(W)$.

Choose coordinates $x_0,\ldots,x_d$ on $W$ with
$p=[1:0:\cdots:0]$.  On $x_0=1$, a zero value and zero directional
derivatives mean that $F(1,x_1,\ldots,x_d)$ has no terms of degree zero or
one.  By homogeneity this is equivalent to $F\in I(p)^2$.  Therefore
\begin{equation}\label{eq:double-containment}
 I_m\subseteq\bigcap_{p\in Z}I(p)^2.
\end{equation}

Suppose $I_m\neq0$ and choose a nonzero homogeneous $F\in I$ of minimal
positive degree $\alpha\leq m$.  For each $p\in Z$, choose
$G_p\in R_{m-\alpha}$ with $G_p(p)\neq0$.  Then $G_pF\in I_m$, so
\eqref{eq:double-containment} and local invertibility of $G_p$ give
$F\in I(p)^2$ locally at every support.  All first partial derivatives of
$F$ consequently vanish on $Z$ and belong to the radical ideal $I$.

If at least one partial derivative is nonzero, it has degree $\alpha-1$ and
contradicts the minimality of $\alpha$.  Otherwise the characteristic is
$p>0$ and every monomial of $F$ has each exponent divisible by $p$.  Because
$\kk$ is algebraically closed, hence perfect, its coefficients have $p$th
roots and
\[
 F=G^p
\]
for a homogeneous polynomial $G$ of degree $\alpha/p<\alpha$.  Radicality of
$I$ and $F\in I$ imply $G\in I$, again a contradiction.  Hence $I_m=0$.

The evaluation map $R_m\to\kk^Z$ is injective and has rank
$\binom{d+m}{d}$.  Every form on $W$ extends to $V$, so this value matrix is
a submatrix of restriction for $H^m$ and its rank is at most both
$\dim S_Z$ and $|Z|$.
\end{proof}

\begin{remark}
The reducedness of $Z$, equivalently radicality of $I(Z)$ in this setting,
is essential to this proof in the Frobenius case.  The argument is elementary but is
consistent with the differential characterization of symbolic powers over
perfect fields; see \cite{DeStefaniGrifoJeffries2020} for a broader modern
framework.
\end{remark}

\section{Proofs of the rank bounds}

\begin{proof}[Proof of Theorem~\ref{thm:main}]
Starting from any support, repeated application of
Lemma~\ref{lem:support-alternative}(i) selects $d+1$ points with linearly
independent representatives.  Lemma~\ref{lem:projective-blocks} makes their
first neighborhoods independent.  Their dual tangent blocks, each of
dimension $d+1$, lie in $S_Z$, so
\[
 \dim S_Z\geq(d+1)^2,
 \qquad |Z|\geq\dim S_Z.
\]

Apply Lemma~\ref{lem:support-alternative}(ii).  If a further point escapes,
the resulting $d+2$ supports are projectively independent and
Lemma~\ref{lem:projective-blocks} gives alternative~(a).  If no point escapes,
the same lemma gives a common $W$ and common Gauss fiber; then
Lemma~\ref{lem:gauss-unisolvence} gives alternative~(b).  Taking the smaller
branch bound proves \eqref{eq:two-branch}.
\end{proof}

\begin{proof}[Proof of Theorem~\ref{thm:growing}]
Write first $m=2q+1$.  If $|Z|\leq q$, part~(i) of
Lemma~\ref{lem:arbitrary-blocks}, followed by degree raising, makes $2Z$
independent in degree $m$.  Its dual tangent span would have dimension
$|Z|(d+1)$ but would lie by absorption in $S_Z$, whose dimension is at most
$|Z|$, a contradiction.  Hence $Z$ contains $q+1$ supports.  Their first
neighborhoods are independent in $H^{2(q+1)-1}=H^m$, so
\[
 \dim S_Z,\ |Z|\geq(q+1)(d+1)=J(d,m).
\]

If $m=2q$, the same argument excludes $|Z|\leq q$.  Choose $q$ supports and
one further support.  Lemma~\ref{lem:arbitrary-blocks}(ii) makes the first
$q$ double blocks and the last simple value independent in degree $2q=m$.
All their dual spaces lie in $S_Z$, giving
\[
 \dim S_Z,\ |Z|\geq q(d+1)+1=J(d,m).
\]
The threshold comparison is immediate in the two parity cases.
\end{proof}

\section{Injective Gauss maps from factorization}

\begin{proposition}[Point-separating factors]\label{prop:factorized-gauss}
Let $H=A\otimes B$ be very ample, and suppose $A$ and $B$ each separate every
ordered pair of distinct closed points.  Then the Gauss map of the complete
$H$-embedding is injective on closed points.
\end{proposition}

\begin{proof}
Fix $p\neq q$.  Choose
\[
 a\in H^0(X,A),\quad a(p)=0,\ a(q)\neq0,
 \qquad
 b\in H^0(X,B),\quad b(p)=0,\ b(q)\neq0.
\]
The product $s=ab\in H^0(X,H)$ belongs to
$H^0(X,H\otimes\mathfrak m_p^2)$ but satisfies $s(q)\neq0$.  Thus the
hyperplane defined by $s$ contains the affine tangent space at $p$ but not
the evaluation line at $q$, which lies in the tangent space at $q$.
Consequently $\widehat T_pX\neq\widehat T_qX$.
\end{proof}

\begin{proof}[Proof of Corollary~\ref{cor:factorized}]
Proposition~\ref{prop:factorized-gauss} makes every Gauss fiber a singleton.
Alternative~(b) of Theorem~\ref{thm:main} would require
$|Z|\geq\binom{d+m}{d}>1$ inside one fiber, so it is impossible.  If
$H=L^r$ with $r\geq2$, use the point-separating factors $L$ and $L^{r-1}$.
\end{proof}

\section{Smooth hypersurfaces realizing the exceptional bound}

We now work over $\CC$.  We first record two elementary interpolation facts.

\begin{lemma}[Simplex-lattice unisolvence]\label{lem:simplex}
For $n,r\geq0$, the lattice
\[
 \Delta_{n,r}=\{(a_1,\ldots,a_n)\in\mathbf Z_{\geq0}^n:
                         a_1+\cdots+a_n\leq r\}
\]
is unisolvent for polynomials of total degree at most $r$.
\end{lemma}

\begin{proof}
If $n=0$ or $r=0$, the assertion is the constant-polynomial case.  Now induct
on $n+r$.  If $f$ vanishes on $\Delta_{n,r}$, its restriction to
$x_n=0$ vanishes on $\Delta_{n-1,r}$, so $f=x_ng$ by induction.  Then
$g(x_1,\ldots,x_{n-1},t+1)$ vanishes on $\Delta_{n,r-1}$ and is zero by
induction.  Thus $f=0$.  The number of lattice points is
$\binom{n+r}{n}$, the dimension of the polynomial space.
\end{proof}

\begin{lemma}[Prescribed triple jets]\label{lem:triple-jets}
Let $p_1,\ldots,p_N$ be distinct points of $\PP^d$ and set $e=3N$.  Then
\[
 H^0(\PP^d,\cO(e))\longrightarrow
 \bigoplus_{i=1}^N\cO_{\PP^d,p_i}/\mathfrak m_{p_i}^3
\]
is surjective.  Moreover, $H^0(\PP^d,\mathcal I_{3Z}(e))$ is basepoint-free
on $\PP^d\setminus Z$.
\end{lemma}

\begin{proof}
For fixed $i$, choose for each $j\neq i$ a linear form vanishing at $p_j$
and nonzero at $p_i$.  The product of their cubes has degree $3(N-1)$, is a
unit at $p_i$, and vanishes to order at least three at every other support.
In affine coordinates centered at $p_i$, the classes of monomials of total
degree at most two form a basis of the local algebra modulo
$\mathfrak m_{p_i}^3$; they are restrictions of global quadrics.  After
multiplying by one further linear form nonzero at every support and adjusting
the local datum by the resulting unit, degree-$e$ forms realize arbitrary
triple jets at $p_i$ and zero jets elsewhere.  Summing over $i$ proves
surjectivity.

For $q\notin Z$, choose a linear form $\ell_i$ vanishing at $p_i$ and
nonzero at $q$ for every $i$.  The product $\prod_i\ell_i^3$ has degree $e$,
belongs to $I_{3Z}$, and is nonzero at $q$.
\end{proof}

\begin{proposition}[A large special Gauss fiber]\label{prop:hypersurface}
For every $d\geq2$ and $N\geq1$ there is a smooth integral hypersurface
$X^d\subset\PP^{d+1}_{\CC}$ of degree $3N$ and a hyperplane
$\Lambda\cong\PP^d$ such that one fiber of the Gauss map contains any
prescribed set of $N$ distinct points of $\Lambda$.
\end{proposition}

\begin{proof}
Use coordinates with $\Lambda=(x_0=0)$ and let
$Z=\{p_1,\ldots,p_N\}\subset\Lambda$.  By
Lemma~\ref{lem:triple-jets}, choose a degree-$e$ form $f_0$ on $\Lambda$
whose constant and linear terms vanish at each $p_i$ and whose quadratic term
there is nondegenerate.  Every member of
\[
 f_0+H^0(\Lambda,\mathcal I_{3Z}(e))
\]
has the same ordinary-double-point jets.  The translation space is
basepoint-free away from $Z$.  The corresponding projective system has base
locus exactly $Z$; its nonempty Bertini open meets the affine chart where the
coefficient of $f_0$ is one.  Hence Bertini gives a member $f$ whose hypersurface
$Y=(f=0)\subset\Lambda$ has ordinary double points at the $p_i$ and is smooth
away from them \cite[III, Corollary~10.9]{Hartshorne1977}.

Consider ambient degree-$e$ forms
\[
 F=f+x_0G,
 \qquad G\in H^0(\PP^{d+1},\cO(e-1)).
\]
The projective linear system spanned by $f$ and
$x_0H^0(\PP^{d+1},\cO(e-1))$ has base locus exactly $Y$.  Its Bertini-smooth
locus outside $Y$ is a nonempty open subset of the irreducible parameter
space, and it meets the affine chart in which the coefficient of $f$ is one.
At a smooth point of $Y$, the differential
of $f$ in a direction tangent to $\Lambda$ is nonzero.  At a marked node,
choose $G(p_i)\neq0$; then
\[
 dF_{p_i}=G(p_i)\,dx_0,
\]
so the ambient hypersurface is smooth there and has tangent hyperplane
$\Lambda$.  Away from $Y$, Bertini makes a general member smooth
\cite[III, Corollary~10.9]{Hartshorne1977}.  The finitely many conditions
$G(p_i)\neq0$ are nonempty open conditions in the same irreducible affine
chart, so their intersection with the Bertini open is nonempty.  The
resulting hypersurface is reduced by smoothness.  If it had two
positive-degree components, they would meet in $\PP^{d+1}$ because $d\geq2$
\cite[I, Theorem~7.2]{Hartshorne1977},
making their union singular; hence it is irreducible and therefore integral.
All marked points lie in the Gauss fiber over $\Lambda$.
\end{proof}

\begin{proof}[Proof of Theorem~\ref{thm:sharp-gauss}]
Take $N=\binom{d+m}{d}$ and identify the affine chart of
$\Lambda\cong\PP^d$ with $\CC^d$.  Let $Z$ be the homogenization of
$\Delta_{d,m}$.  Lemma~\ref{lem:simplex} says that no nonzero degree-$m$
form on $\Lambda$ vanishes on all of $Z$.  Equivalently, their degree-$m$
evaluation vectors span the full vector space $\Sym^m(\widehat\Lambda)$ of
dimension $N$.  Here we use the canonical dual identification
$H^0(\Lambda,\cO_\Lambda(m))^*\cong\Sym^m(\widehat\Lambda)$.

Apply Proposition~\ref{prop:hypersurface} to these supports, obtaining a
smooth degree-$e=3N$ hypersurface $X$.  Since $e>m$ and $d+1\geq3$, the
restriction sequence
\[
 0\longrightarrow\cO_{\PP^{d+1}}(m-e)
 \longrightarrow\cO_{\PP^{d+1}}(m)
 \longrightarrow\cO_X(m)\longrightarrow0
\]
and the vanishing of intermediate cohomology of projective space
\cite[III, Theorem~5.1]{Hartshorne1977} show that
\[
 H^0(\PP^{d+1},\cO(m))\xrightarrow{\sim}H^0(X,\cO_X(m)).
\]
Thus the complete $\cO_X(m)$-embedding is the restriction of the ambient
$m$th Veronese embedding.

If $p=[v]\in Z$, Proposition~\ref{prop:hypersurface} gives
$\widehat T_pX=\widehat\Lambda$.  The affine tangent space after the Veronese
embedding is
\[
 v^{m-1}\widehat\Lambda\subseteq\Sym^m(\widehat\Lambda).
\]
The differential carries the scalar factor $m$, which is nonzero and
irrelevant to this subspace because the construction is over $\CC$.
The right side is exactly the span of the evaluation vectors of $Z$.
Therefore every tangent block lies in $S_Z$ and $Z$ is absorbing.  The same
unisolvence gives $\dim S_Z=N=|Z|$.
\end{proof}

\begin{remark}
For $(d,m)=(2,3)$ the construction gives ten absorbing points on a smooth
degree-$30$ surface in $\PP^3$, so the exact value ten is not confined to a
linearly embedded plane.  In higher parameters the example is sharp for the
exceptional branch but need not attain the smaller global branch bound.
\end{remark}

\section{Relation to adjacent notions and limitations}

The equality of value and first-jet kernels in
Lemma~\ref{lem:kernel} is stronger than an ordinary Terracini defect.  It says
that the full double-point rank has collapsed to the value rank.  Conversely,
a Terracini defect alone does not imply absorption
\cite{BallicoChiantini2021,GaluppiEtAl2023}.

Under absorption one also has
\[
 S_Z=\Span\{\widehat T_{\varphi_m(p)}\varphi_m(X):p\in Z\}.
\]
If $A\subseteq Z$ is minimal with the same tangent span, then
$Z\setminus A$ lies in the strong base locus of $A$.  This places the problem
inside the framework of \cite{BallicoBrambillaSantarsiero2026}, while the
additional equality with a span of evaluations is the special constraint
used here.  Strong base loci and Terracini loci are not comparable in general
\cite{BallicoBrambillaSantarsiero2026}.
For classical background on Gauss maps, tangent loci, and dual varieties, see
\cite{Zak1993}.
Multitangent hyperplanes also have a substantial enumerative and dual-variety
literature.  Vainsencher counts $n$-fold tangent hyperplanes to surfaces in
specific generic settings for $1\leq n\leq6$ \cite{Vainsencher1995}, while
Holweck studies bitangent components in singular loci of dual varieties
\cite{Holweck2011}.  Those works concern enumerative or dual-singularity
questions; they do not by themselves establish the value-span absorption
constraint used here or priority for the present construction.

The following limitations are material.
\begin{itemize}
\item Theorem~\ref{thm:main} is stated over algebraically closed fields.  A
version over a nonclosed perfect field would require a careful formulation for
closed points and residue fields; it is not asserted here.
\item The hypersurface construction is over $\CC$.  It uses ordinary double
points with nondegenerate quadratic parts and a standard Bertini argument;
no positive-characteristic version is claimed.
\item The construction proves existence, not an explicit globally smooth
equation.  Finite rank replays can certify interpolation and absorption
matrices, but cannot replace Bertini or prove the universal theorem.
\item Neither equality configurations nor varieties with large special Gauss
fibers are classified.  The simplex lattice is one unisolvent fixture among
many.
\item The cited literature comparison is selective.  The paper claims the
displayed deductions and constructions, not exhaustive priority for tangent
containment, jet separation, or special Gauss fibers.
\end{itemize}

\section{Conclusion}

Point-span tangent absorption admits a uniform characteristic-free structure:
support escape gives independent tangent blocks, while failure of escape puts
all supports into a single original Gauss fiber and forces degree-$m$
unisolvence there.  Arbitrary-support interpolation adds the growing bound
$J(d,m)$.  Factorization eliminates the Gauss branch when the original
$H$-embedding is complete, whereas smooth complex hypersurfaces with prescribed nodal hyperplane
sections realize its binomial estimate exactly.  These results sharpen the
rank picture in arbitrary characteristic without asserting a classification
or relying on computational fixtures as proof.  In characteristic zero the
stronger exact global floor is supplied by the companion work
\cite{ErikssonExactARR2026}.

\begin{thebibliography}{99}

\bibitem{BallicoBrambillaSantarsiero2026}
E.~Ballico, M.~C.~Brambilla, and P.~Santarsiero,
\emph{On the strong base locus of a projective variety},
arXiv:2603.15103 (2026).
\url{https://arxiv.org/abs/2603.15103}

\bibitem{BallicoChiantini2021}
E.~Ballico and L.~Chiantini,
\emph{On the Terracini locus of projective varieties},
Milan J. Math. \textbf{89} (2021), 1--17.
\url{https://doi.org/10.1007/s00032-020-00324-5}

\bibitem{BeltramettiDiRoccoSommese1999}
M.~C.~Beltrametti, S.~Di Rocco, and A.~J.~Sommese,
\emph{On generation of jets for vector bundles},
Rev. Mat. Complut. \textbf{12} (1999), 27--45.
\url{https://doi.org/10.5209/rev_REMA.1999.v12.n1.17182}

\bibitem{DeStefaniGrifoJeffries2020}
A.~De Stefani, E.~Grifo, and J.~Jeffries,
\emph{A Zariski--Nagata theorem for smooth $\mathbf Z$-algebras},
J. Reine Angew. Math. \textbf{761} (2020), 123--140.
\url{https://doi.org/10.1515/crelle-2018-0012}

\bibitem{ErikssonB219}
L.~Eriksson,
\emph{B219---Special Gauss fibers can be arbitrarily large},
public working note (2026),
\href{https://github.com/lluiseriksson/hodge-conjecture-research-front/blob/main/proofs/B219-arbitrarily-large-special-gauss-fibers.md}{repository record B219}.

\bibitem{ErikssonB220}
L.~Eriksson,
\emph{B220---A factorized polarization has injective Gauss map},
public working note (2026),
\href{https://github.com/lluiseriksson/hodge-conjecture-research-front/blob/main/proofs/B220-factorized-polarization-has-injective-gauss-map.md}{repository record B220}.

\bibitem{ErikssonExactARR2026}
L.~Eriksson,
\emph{The Exact Universal Rank Floor for Point-Span Tangent Absorption:
Projection, Fattening, and Common-Tangent Extremizers},
ARR-2026-2MHNZRRJP49Y9SWP, v1 (2026).
\url{https://arr-research.github.io/papers/ARR-2026-2MHNZRRJP49Y9SWP/}

\bibitem{GaluppiEtAl2023}
F.~Galuppi, P.~Santarsiero, D.~A.~Torrance, and E.~T.~Turatti,
\emph{Geometry of first nonempty Terracini loci},
arXiv:2311.09067 (2023).
\url{https://arxiv.org/abs/2311.09067}

\bibitem{Hartshorne1977}
R.~Hartshorne,
\emph{Algebraic Geometry},
Graduate Texts in Mathematics, vol.~52, Springer, New York, 1977.

\bibitem{Holweck2011}
F.~Holweck,
\emph{Singularities of duals of Grassmannians},
J. Algebra \textbf{337} (2011), 369--384.
\url{https://doi.org/10.1016/j.jalgebra.2011.04.023}

\bibitem{Vainsencher1995}
I.~Vainsencher,
\emph{Enumeration of $n$-fold tangent hyperplanes to a surface},
J. Algebraic Geom. \textbf{4} (1995), 503--526.
\url{https://arxiv.org/abs/alg-geom/9312012}

\bibitem{Zak1993}
F.~L.~Zak,
\emph{Tangents and Secants of Algebraic Varieties},
Translations of Mathematical Monographs, vol.~127,
American Mathematical Society, Providence, 1993.

\end{thebibliography}

\end{document}
